\chapter{Topological Spaces}
\pagenumbering{arabic}

\section{Topological Spaces}
\begin{definition}\label{def:topology}
    A \textui{topology} on a set $X$ is a collection $T$ of subsets of $X$ such that:
        \begin{enumerate}[label = (\arabic*),itemsep=1pt,topsep=3pt]
            \item $\emptyset, X \in T$;
            \item if $\{U_i\}_{i \in I} \subseteq T$, then $\bigcup_{i \in I}U_i \in T$;
            \item if $\{U_i\}_{i = 1}^n \subseteq T$, then $\bigcap_{i = 1}^n \in T$.
        \end{enumerate}
    We call the pair $(X,T)$ a \textui{topological space}.
\end{definition}

\begin{example}
    Let $X = \{a,b,c\}$. Then:$T_1 := \{\emptyset, \{a,b,c\},\{b\},\{a,b\},\{b,c\}\}$ and $T_2 := \{\emptyset,\{b\},\{a,b,c\}\}$ are topologies on $X$. However, $T_3 =\{\emptyset, \{a\},\{b\},\{a,b,c\}\}$ is \textit{not} a topology on $X.$
\end{example}

\begin{example}[Every set admits a topology]
    Let $X$ be a set. The \textit{discrete topology} on $X$ is $T = \cP(X)$. The \textit{trivial topology} on $X$ is $T = \{\emptyset, X\}$.
\end{example}

\begin{definition}
    Let $T_1$ and $T_2$ be topologies on a set $X$. If $T_1 \subseteq T_2$, then we say $T_2$ is \textui{finer} than $T_1$, and likewise that $T_1$ is \textui{coarser} than $T_2$.
\end{definition}

\begin{example}\label{example:metric-on-Rn}
    Let $\vec{x} = (x_1,...,x_n)$ and $\vec{y} = (y_1,...,y_n)$ be vectors in $\bfR^n$. Define $d(\vec{x},\vec{y}) := \sqrt{(x_1-y_1)^2 + ... + (x_n - y_n)^2}$. This is a metric on $\bfR^n$, and we will see this induces a topology on $\bfR^n$.
\end{example}

\begin{definition}
    Let $(X,d)$ be a metric space, $x \in X$, and $\epsilon > 0$. The \textui{open ball centered at $x$} \textui{with radius $\epsilon$} is denoted $B(x,\epsilon) = \{y \in X \mid d(x,y) < \epsilon\}$.
\end{definition}

\begin{definition}
    Let $(X,d)$ be a metric space. A set $U \subseteq X$ is \textui{open} if for all $x \in U$, there exists $\epsilon >0$ such that $B(x,\epsilon) \subseteq U$. The collection of open sets of $X$ is denoted $\tau_X$.
\end{definition}

\begin{proposition}
    Let $(X,d)$ be a metric space. Then $(X,\tau_X)$ is a topological space.
\end{proposition}
    \begin{proof}
        Clearly $\emptyset,X \in \tau_X$. Let $\{U_i\}_{i \in I} \subseteq X$ be an arbitrary family of open sets. Let $x \in \bigcup_{i \in i}U_i$. Then $x \in U_i$ for some $i$. So there exists an $\epsilon > 0$ such that $B(x,\epsilon) \subseteq U_i \subseteq \bigcup_{i \in I}U_i$. Now let $\{U_i\}_{i = 1}^n$ be a finite collection of open sets. Let $x \in \bigcap_{i = 1}^n U_i$. Then $x \in U_i$ for all $i = 1,...,n$. So there exists $\epsilon_i > 0$ such that $B(x,\epsilon_i) \subseteq U_i$ for all $i = 1,...,n$. Define $\epsilon := \min\{\epsilon_1,...,\epsilon_n\}$. Then $B(x,\epsilon) \subseteq U_i$ for all $i=1,...,n$. Hence $B(x,\epsilon) \subseteq \bigcap_{i = 1}^n U_i$. Thus $(X,\tau_X)$ is a topological space.
    \end{proof}

\begin{example}
    Using the same terminology from Example~\ref{example:metric-on-Rn}, since $d(\vec{x},\vec{y})$ is a metric, $(\bfR^n, \tau_{\bfR^n})$ is a topological space.
\end{example}

\section{Topologies Generated by a Basis}
\begin{definition}
    Let $X$ be a set. A \textui{basis for a topology on $X$} is a collection $\cB$ of subsets of $X$ such that:
        \begin{enumerate}[label = (\arabic*),itemsep=1pt,topsep=3pt]
            \item for all $x \in X$, there exists $B \in \cB$ such that $x \in \cB$.
            \item If $x \in B_1 \cap B_2$, there exists $B_3 \in \cB$ such that $x \in B_3 \subseteq B_1 \cap B_2$.
        \end{enumerate}
\end{definition}

\begin{remark}
    Let $(X,d)$ be a metric space and $\cB$ a basis for a topology on $X$. It follows from the above definition that for any $U \in \tau_X$, there exists an open basis element $B \in \cB$ such that $x \in B \subseteq U$.
\end{remark}

\begin{example}\label{example:basis-1}
    The sets $\cB_1 := \{(a,b) \mid a<b\}$ and $\cB_2 := \{(a,b) \mid a < b, a,b \in \bfQ\}$ are bases for a topology on $\bfR$. These definitions extend nicely to $\bfR^n$ by considering open balls.
\end{example}

\begin{example}[Every set admits a basis for a topology]\label{example:basis-2}
    Let $X$ be a set. Then $\cB := \{\{p\} \mid p \in X\}$ is a basis for a topology on $X$.
\end{example}

\begin{proposition}\label{prop:top-gen-by-basis}
    Let $X$ be a set and $\cB$ a basis for a topology on $X$. Let $T$ be the set containing the empty set and every subset of $X$ which is equal to a union of basis elements of $\cB$. Then $(X,T)$ is a topological space.
\end{proposition}
    {\color{red} \begin{proof}
        
    \end{proof}}

\begin{definition}
    Let $X$ be a set and $\cB$ a topology on $X$. The topology defined in Proposition~\ref{prop:top-gen-by-basis} is called the \textui{topology generated by $\cB$}.
\end{definition}

\begin{example}
    From Example~\ref{example:basis-1}, $\cB_1$ and $\cB_2$ generate the standard topology on $\bfR^n$. From Example~\ref{example:basis-2}, $\cB$ generates the discrete topology on $X$.
\end{example}

\section{Product Topologies}

\begin{definition}
    Let $(X,T_1)$ and $(Y,T_2)$ be topological spaces. The \textui{product topology on $X \times Y$} is the topology generated by the basis $\cB := \{U \times V \mid U \in T_1, V \in T_2\}$.
\end{definition}

\begin{remark}
    Note that $\cB$ is not necessarily a topology on $X \times Y$. Indeed, 
    consider $X = Y = \{0,1\}$ in the discrete topology. Then the sets $\{0\}\times\{0\}$ and $\{0\}\times\{1\}$ are in $\cB$ as per the above definition, but their union $\{(0,0),(0,1)\}$ is not. Hence $\cB$ is not a topology.
\end{remark}

\begin{example}\label{example:product-bases}
    Recall from Example~\ref{example:basis-1} that $\cB := \{(a,b) \mid a < b\}$ generates the standard topology on $\bfR$. Consider the following to bases for a topology on $\bfR^2$:
        \begin{equation*}
        \begin{split}
            \cB_1 &:= \{B(x,r) \mid x \in \bfR^2, r > 0\}, \\
            \cB_2 &:= \{I_1 \times I_2 \mid I_1, I_2 \in \cB\}.
        \end{split}
        \end{equation*}
    The basis elements of $\cB_1$ are open balls of radius $r$, whilst the basis elements of $\cB_2$ are open squares. Note that the basis $\cB_2$ interprets $\bfR^2$ as the product $\bfR \times \bfR$.

    This idea can be further generalized to $\bfR^3$. Consider:
        \begin{equation*}
        \begin{split}
            \cB_3 &:= \{B(x,r) \mid x \in \bfR^3 , r > 0\} \\
            \cB_4 &:= \{I_1 \times I_2 \times I_3 \mid I_1,I_2,I_3 \in \cB\} \\
            \cB_5 &:= \{ B \times I \mid B \in \cB_1, I \in \cB\}.
        \end{split}
        \end{equation*}
    The basis $\cB_4$ interprets $\bfR^3$ as $\bfR \times \bfR \times \bfR$, whilst the basis $\cB_5$ interprets $\bfR^3$ as $\bfR^2 \times \bfR$.
\end{example}

\begin{proposition}
    The bases $\cB_1$ and $\cB_2$ from Example~\ref{example:product-bases} generate the same topologies on $\bfR^2$.
\end{proposition}
    {\color{red} \begin{proof}
        
    \end{proof}}

\section{Subspace Topology}

\begin{proposition}\label{prop:subspace-top}
    Let $(X,T)$ be a topological space. If $Y \subseteq X$, then $T_Y := \{Y \cap U \mid U \in T_X\}$ is a topology on $Y$.
\end{proposition}
    {\color{red} \begin{proof}
        
    \end{proof}}

\begin{definition}
    Let $(X,T)$ be a topological space and $Y \subseteq X$. The \textui{subspace topology on $Y$} is the topology defined in Proposition~\ref{prop:subspace-top}. We call $Y$ a \textui{subspace of $X$} if it is paired with the subspace topology.
\end{definition}

\begin{example}
    Consider the interval $[2,4)$ with the subspace topology inherited from $\bfR$. Note that $(3,4)$ is open in $[2,4)$ and $\bfR$, while $[2,3)$ is open in $[2,4)$ but \textit{not} open in $\bfR$.
\end{example}

\begin{example}
    Consider the integers $\bfZ$ with the subspace topology inherited from $\bfR$. For any $n \in \bfZ$, we have $\{n\} = \bfZ \cap \left( n - \frac{1}{2}, n+ \frac{1}{2} \right)$. Since $\left( n - \frac{1}{2}, n+ \frac{1}{2} \right)$ is open in $\bfR$, each singleton in $\bfZ$ is open. Hence the subspace topology on $\bfZ$ inherited from $\bfR$ is the discrete topology.
\end{example}

\begin{proposition}
    If $\cB$ is a basis for a topology on a set $X$, and $Y$ is a subset of $X$, then $\cB_y := \{Y \cap U \mid U \in \cB\}$ is a basis for $Y$.
\end{proposition}
    {\color{red} \begin{proof}
        
    \end{proof}}

\begin{example}
    Recall that $\cB := \{(a,b) \mid a < b\}$ is a basis for the standard topology on $\bfR$. Consider $\cB_\bfQ := \{(a,b) \cap \bfQ \mid a < b\}$. Since $\bfQ$ is dense in $\bfR$, we have that $(a,b) \cap \bfQ \neq \emptyset$. In fact, each nonempty open set of $\bfQ$ with this particular subspace topology is infinite, hence it is not the discrete topology.
\end{example}

\section{Closed Sets and Limit Points}

\begin{definition}
    A subset $A$ of a topological space $X$ is \textui{closed in $X$} if $A^c$ is open.
\end{definition}

\begin{example}
    Let $(X,T)$ be a topological space. Then $\emptyset$ and $X$ are closed.
\end{example}

\begin{example}\label{example:singletons-closed-in-Rn}
    If $x \in \bfR^n$, then $\{x\}$ is closed in $\bfR^n$. Indeed, for any $p \in \bfR^n \setminus \{x\}$, we have that $p \in B \left( p , \frac{\lnorm p - x \rnorm}{3} \right) \subseteq \bfR^n$.
\end{example}

\begin{example}
    Given any $x \in \bfR^n$, the set $\{y \in \bfR^n \mid \lnorm x - y \rnorm \leq \epsilon\}$ is closed. The proof of this fact follows identically to Example~\ref{example:singletons-closed-in-Rn}.
\end{example}

\begin{example}
    Consider the topological space $(X,\cP(X))$. Since every set is open, every set must be closed.
\end{example}

\begin{proposition}
    Let $X$ be a topological space. Then:
    \begin{enumerate}[label = (\arabic*),itemsep=1pt,topsep=3pt]
        \item $\emptyset$ and $X$ are closed;
        \item If $\{C_i\}_{i \in I}$ is an arbitrary family of closed sets, then $\cap_{i \in I}C_i$ is closed;
        \item If $\{C_i\}_{i =1}^n$ is a finite collection of closed sets, then $\bigcup_{i = 1}^n C_i$ is closed.
    \end{enumerate}
\end{proposition}
    \begin{proof}
        This follows from De Morgan's laws and Definition~\ref{def:topology}.
    \end{proof}

\begin{example}
    Every finite subset of $\bfR^n$ is closed. Indeed, any finite subset can be written as the finite union of singletons, which are closed in $\bfR^n$ by Example~\ref{example:singletons-closed-in-Rn}.
\end{example}

\begin{proposition}
    Let $Y$ be a subspace of a topological space $X$. Then a set $A \subseteq Y$ is closed in $Y$ if and only if $A = Y \cap C$, where $C$ is a closed set in $X$.
\end{proposition}
    {\color{red} \begin{proof}
        
    \end{proof}}

\section{Closure and Interior}

\begin{definition}
    Let $(X,T)$ be a topological space and $A\subseteq X$.
    \begin{enumerate}[label = (\arabic*),itemsep=1pt,topsep=3pt]
        \item The \textui{interior of $A$} is $A^o := \bigcup \{U \mid U \subseteq A, U \text{ open }\}$.
        \item The \textui{closure of $A$} is $\overline{A} := \bigcap\{C \mid C \supseteq A, C \text{ closed }\}$.
    \end{enumerate}
\end{definition}

\begin{remark}
\phantom{a}
\begin{enumerate}[label = (\arabic*),itemsep=1pt,topsep=3pt]
    \item Clearly $A^o \subseteq A \subseteq \overline{A}$. The interior of any set is open and is the largest open subset containing it. The closure of any set is closed and is the smallest closed subset contained in it. A set is open if and only if it is equal to its interior. A set is closed if and only if it is equal to its closure. If $C$ is any closed set in $X$ and $A \subseteq C$, then $\overline{A} \subseteq C$.
\end{enumerate}
\end{remark}

\begin{example}
    If $F \subseteq \bfR^n$ is finite then $F^o = \emptyset$ and $\overline{F} = F$.
\end{example}

\begin{example}
    Consider the interval $(2,3]$ in $\bfR$. Then $(2,3]^o = (2,3)$ and $\overline{(2,3]} = [2,3]$.
\end{example}

\begin{example}
    In the discrete topology, every set is open and closed, hence every set equals its interior and closure.
\end{example}

\begin{example}
    Consider the set $A:= \{1,\frac{1}{2}, \frac{1}{3},...\}$ as a subset of $\bfR$. Then $A^o = \emptyset$ and $\overline{A} = A \cup \{0\}$.
\end{example}

\begin{proposition}
    Let $(X,T)$ be a topological space and $A \subseteq X$.
    \phantom{a}
    \begin{enumerate}[label = (\arabic*),itemsep=1pt,topsep=3pt]
        \item $p \in \overline{A}$ if and only if for all open sets $U$ in $X$ containing $p$, we have $U \cap A \neq \emptyset$.
        \item If $\cB$ is a basis for $T$, then $p \in \overline{A}$ if and only if for all basis elements $B \in \cB$ containing $p$, we have $B \cap A \neq \emptyset$.
    \end{enumerate}
\end{proposition}
    {\color{red} \begin{proof}
        
    \end{proof}}

\section{Limit Points of a Set}

\begin{definition}
    Let $(X,T)$ be a topological space. Let $A \subseteq X$. A point $p \in A$ is a \textui{limit point of $A$} if, for all open $U \subseteq X$ containing $p$, then $U \cap (A\setminus\{p\}) \neq \emptyset$.
\end{definition}

